\section{Results}
\label{sec:results}

 
\subsection*{Pre-trial phase}

Table \ref{tab:twfe_table} provides the TWFE results for all the variables in the analysis, and should serve as a benchmark for comparisons. Figure \ref{fig:att_first} reports the estimated impact of federal prison construction near a municipality on the number of cases processed, individuals sent to pre-trial detention, and individuals released after the first hearing with the control judge. The estimated effects are positive and statistically significant across all three outcomes. The overall ATT—defined as the average effect of treatment across all treated groups and all post-treatment periods, weighted by group size—indicates that municipalities where a federal prison was built experienced an approximate 3.8\% increase (CSDID with controls) in processed cases, pre-trial detentions, and releases relative to the control group. Given the similarity of coefficients across these outcomes, the evidence suggests that the results are driven by a general increase in case loads, consistent with the interpretation that prison construction enhanced the state government’s ability or willingness to prosecute.

Although this measure is imperfect—since zero values in the denominator require imputation and adding a dummy for the extensive margin may constitute a bad control—the results for ratios provide additional insight. The share of individuals sent to pre-trial detention out of all processed cases increases significantly, while the share of individuals released remains statistically indistinguishable from zero.\footnote{See Section \ref{sec:robust_log} for a discussion on the possible biases introduced from using rates as a dependent variable, as well as possible solutions not explored in this paper.} This pattern suggests that prison construction may have altered judicial behavior, making judges more likely to impose pre-trial detention when additional capacity becomes available.

Figure \ref{fig:event_first_total} presents the event-study estimates for the number of processed individuals (top left), individuals sent to pre-trial detention (top right), and individuals released (bottom). Across the three outcomes, there is no evidence of systematic pre-treatment trends, although some coefficients are sporadically significant. Importantly, the ATTs mask considerable temporal heterogeneity: post-treatment effects follow an inverted U-shape, increasing in the medium term after prison construction and gradually declining over the longer run (within ten years). Finally, Figure \ref{fig:event_first_rate} shows event study estimates for the rates, which show no differential pre-trends and confirm the results of the ATTs.

\subsection*{Sentencing outcomes (first instance courts)}

Figure \ref{fig:att_sent} reports the estimated impact of federal prison construction near a municipality on judicial sentencing outcomes: the number of individuals sentenced, found guilty and sent to prison, found guilty and ordered to pay a fine, and acquitted (panels 1–4, top to bottom). These outcomes correspond to the final decision of a judge after each individual has completed the judicial process. The results present a consistent pattern. The overall number of sentenced individuals rises significantly, driven primarily by a 4\% increase in guilty verdicts resulting in prison sentences. By contrast, the effects on fines and acquittals are precisely estimated as null. This suggests that the results reflect changes in judicial behavior rather than spillovers from the increased caseloads documented above.

Figure \ref{fig:event_sentence_total} presents event-study estimates for these variables. There is no evidence of pre-treatment differentials, lending credibility to the parallel trends assumption. Post-treatment dynamics are similar to those in the processed-cases analysis: an initial increase followed by a pronounced decline in later periods. The similarity of these dynamic patterns, along with the sharp drop at longer horizons, reveals an intriguing trend explored further in Section \ref{sec:robust_log}.

Turning to ratios, ATT estimates show no significant effects for the average length of sentences, nor for the ratios of guilty-to-sentenced, prison-to-guilty, or fine-to-guilty. However, event-study estimates in Figure \ref{fig:event_sentence_rates} reveal distinct post-treatment dynamics: municipalities near a new prison exhibit a positive differential trend in most ratios after treatment, with the exception of fines relative to guilty verdicts. 

Finally, two points deserve emphasis. First, TWFE estimates are broadly similar to CSDID results, with only a few exceptions. Nonetheless, as shown later, event-study estimates diverge markedly, likely due to negative weighting in TWFE. For this reason, all interpretations are done using CSDID with controls. Second, the similarity between CSDID estimates with and without controls strengthens the case for the parallel trends assumption: the addition of covariates—intended to improve its plausibility—produces little change in point estimates.

\subsection*{Identifying assumptions}

As discussed in the previous section, identifying the Average Treatment on the Treated (ATT) effects of federal prison construction on local sentencing outcomes does not require assuming that the timing or location of prison construction was random. Instead, identification relies on a weaker parallel trends assumption: in the absence of prison construction, municipalities located near new facilities would have exhibited similar incarceration trends as those farther away. Although this assumption is inherently untestable, I provide substantial empirical evidence supporting its plausibility.

One potential concern is that the executive branch strategically chose the timing and location of federal prison construction in response to pre-existing trends in local prison capacity. Under this alternative explanation, federal prisons would have been built precisely in areas where state-level prisons were becoming increasingly overcrowded, in order to alleviate local capacity pressures. I address this concern in two ways. First, as documented earlier, the 2009 decision to construct private federal prisons was publicly justified by President Felipe Calderón as part of a nationwide effort to relocate individuals convicted of federal crimes into federal facilities, since many were being housed in state prisons. This rationale suggests that the policy was driven by jurisdictional considerations rather than by localized trends in state prison overcrowding. The problem was not uneven local congestion, but rather insufficient federal infrastructure to comply with existing legal mandates.

Second, if prison construction had been systematically targeted toward areas experiencing worsening local overcrowding, we would expect to observe differential pre-treatment trends in overcrowding between treated and control municipalities. Figures \ref{fig:capacity} and \ref{fig:capacity_robust} show no evidence of such pre-trends in relative overcrowding at the local level, providing direct evidence against this selection mechanism.

A related concern is that the executive—particularly under President Felipe Calderón—may have strategically directed prison construction toward regions experiencing rising crime or intensified enforcement. The War on Drugs was publicly launched in Michoacán with the deployment of the army on December 11, 2006. However, Table \ref{tab:app1tab} shows that the federal prison in Michoacán was not completed until 2016. In contrast, the federal prisons built after the start of the war—whose construction likely began during the early years of the conflict—were located in the Southeast (Veracruz and Tabasco), while Michoacán is in the West. This pattern suggests that the crackdowns were not accompanied by contemporaneous increases in federal incarceration capacity in the regions most affected by the initial enforcement efforts. Moreover, if this were the case, we would expect to observe differential pre-treatment trends in criminal justice activity, such as the total number of cases processed by courts. Figure \ref{fig:event_first_total} (top left) shows no such pre-trends, further supporting the validity of the parallel trends assumption.

\subsection*{Mechanisms}

Understanding the mechanisms behind changes in sentencing decisions is essential for interpreting the effects of the presence of bottlenecks in the judicial system. Several channels may be at play. First, if the government expands policing in parallel with new prison capacity, observed increases in sentencing may simply reflect a larger caseload rather than genuine shifts in judicial behavior. As shown in the previous section, both the total number of individuals processed and the total number of sentences—particularly those resulting in prison terms or pre-trial detention—rose in a similar manner.

However, additional evidence suggests that this channel alone cannot fully account for the observed increases. While both pre-trial detentions and releases rose during the pre-trial phase, the rate of pre-trial detention relative to the number of individuals processed increased, whereas the release rate remained stable. Taken together, these results point to two dynamics: a higher overall caseload, but also a likely change in judicial behavior, as judges appeared less constrained to impose pre-trial detention once additional prison capacity became available. This interpretation is reinforced by the fact that both overall sentencing and the number of individuals sent to prison increased significantly, while acquittals and sentences involving fines showed no change.

Now, during periods of overcrowding, judges may have issued more lenient sentences due to limited available space; once new facilities become operational, sentencing could return to prior levels, generating an increase in places where new capacity becomes available. To test this hypothesis, I use the geolocated panel dataset of state and federal prisons and run the following regression: 

\begin{equation}
\label{eq:eq3}
    Y_{pt} = \phi + \beta \, TreatPost_{pt} + \gamma_t + \tau_p + \epsilon_{pt},
\end{equation}

where $Y_{pt}$ is the inverse hyperbolic sine transformation of the relative overcrowding of each state prison $p$ in month $t$ (at the bimonthly level). As before, $\gamma_t$ and $\tau_p$ denote time and prison fixed effects. The key difference from the baseline specification lies in the definition of treatment. Here, $TreatPost_{pt}$ is a dummy that equals one once a federal prison becomes operational within a 400 km radius of a state prison that, at least once during the study period, housed a significant number of federal prisoners. As discussed in the background section, agreements between federal and state authorities allow the transfer of inmates across jurisdictions, but data on such arrangements is unavailable. To account for this, I classify a state prison as having such an agreement if the maximum share of federal prisoners it hosted during the period exceeded the median. This approach captures both the prisons that at some point received federal inmates and those most likely to be structurally affected by increases in federal capacity.  

Figure \ref{fig:capacity} presents the event-study estimates of Specification \ref{eq:eq3} using the method proposed by \cite{Callaway}. The results indicate that the construction of a federal prison leads to a significant reduction in overcrowding in state facilities with federal-receiving capabilities. The event-study dynamics show no evidence of differential pre-trends and reveal a persistent, long-term decline in overcrowding following construction, consistent with a structural expansion of capacity. Within six years of a federal facility’s opening, municipalities located within its catchment area display substantially lower levels of overcrowding. These results are robust to alternative definitions of treatment: using thresholds of 300 km and 500 km produces similar patterns, as shown in Figure \ref{fig:capacity_robust}.

\subsection*{Robustness checks}

\subsubsection*{Robustness to other HTE methods}

Figures \ref{fig:att_first_robust} and \ref{fig:att_sent_robustness} present the main results using alternative estimators: \cite{Borusyak-DIDImp} (didimputation) and \cite{SunAb} (SunAb). Across all dependent variables in levels, the ATT estimates are statistically indistinguishable, with closely aligned point estimates. The \cite{SunAb} estimator is effectively equivalent to \cite{Callaway} without controls, since both rely on 2x2 group comparisons across groups (excluding forbidden comparisons, such as using treated units as controls). The motivation for including an imputation-based method lies in its efficiency under homoskedasticity, as demonstrated by \cite{Borusyak-DIDImp}. Although some degree of serial correlation across units is likely, efficiency gains are still expected even under low correlation \citep{Chaise-ReviewDID}. In line with these theoretical results, imputation-based estimates prove more efficient than the event-by-event HTE-robust methods.

The robustness checks require careful interpretation. In general, results are consistent across specifications, but not uniformly across all variables. For instance, rate-based variables yield differences across estimators: the rate of pretrial detention to total processed displays a sign change and significance only with the imputation method. This likely reflects two issues—first, the need to impute zeros for municipalities without cases and include controls for this, and second, the possibility that the denominator itself is affected by the treatment. 

Moreover, in some cases, efficiency gains from imputation methods transform similar point estimates into statistically significant results. Three variables deserve special attention: the rate of pretrial detention to total processed, the log of individuals sentenced not guilty, and sentencing length. To examine these, I plot event-study estimates using \cite{Gardner} and \cite{Liu-FECT} together with \cite{Borusyak-DIDImp}. For sentencing length, Figure \ref{fig:sent_intensive_event_robust} shows post-treatment patterns that look similar across methods, with a clear increase after prison construction. The imputation methods estimate this more precisely, making the ATT significant. The interpretation therefore remains the same: even if ATT estimates differ across methods, the event studies point to a positive post-treatment trend in sentencing length.

The other two variables offer less consistent results. Figures \ref{fig:event_study_absolutoria} and \ref{fig:event_study_pretrial} show that the different methods give different answers. For non-guilty sentences, both sets of methods show a small increase about two years after construction followed by a decline, but the imputation methods make the increase larger and significant, while the event-by-event estimates do not. For the pretrial detention rate, event-by-event methods suggest a lasting increase, while imputation methods show a negative trend. Because of these differences, it is hard to reach firm conclusions for these two outcomes. 

\subsubsection*{Robustness to transformations of the dependent variable}
\label{sec:robust_log}
As noted in the main results, all models for variables in levels are estimated using a log transformation of the dependent variable. This choice makes treatment effects easier to interpret as percentage changes. It also helps address skewness in the data: the number of individuals processed in the judicial system is likely higher in municipalities with larger populations (or other correlated factors), and the log transformation reduces the influence of outliers in the estimation.

However, despite these advantages, \cite{ChenRoth2024} highlight several biases that can arise when using log transformations to estimate treatment effects. These concerns apply not only to ATEs but also to DiD settings (ATTs). The main issue is that in contexts with an extensive margin—where treatment can change an outcome from zero to positive—treatment effects depend on the arbitrary choice of scale. By simply re-scaling the log-transformed dependent variable by a scalar $a$, one could generate any treatment effect. This problem arises because percentage effects are not well defined for observations that move from zero to nonzero after treatment. In such cases, the units of the outcome implicitly determine the weight placed on the extensive margin. In this context, treatment effects are not well defined for municipalities with no individuals processed or sentenced at time $t$ when their outcomes shift from zero to positive—for example, municipalities that would process no cases without a nearby prison but would process at least one once a new facility is built.

I provide evidence that these theoretical concerns do not bias the main results presented in the previous section. Figure \ref{fig:att_sent_robustness_chen} reports the ATTs for all outcomes in levels estimated with the $log(y+1)$ transformation. Two points are worth highlighting. First, across all variables except one, there is no evidence of an extensive-margin treatment effect ($1\{y>0\}$): ATT estimates are precisely estimated null effects, with the only exception being the not-guilty outcome. Second, results using the hyperbolic sine transformation ($arcsinh(y) = \log(y + \sqrt{1 + y^2})$) closely match those from the log-like specification. Finally, because both $arcsinh(y)$ and $log(y+1)$ suffer from the unit-dependence problem, I also estimate ATTs using the transformation proposed by \cite{ChenRoth2024}, which calibrates the extensive margin by setting $y=0$ to $-x$ (where $x = \min\{y>0\}$) and $\log(y)$ for $y>0$, so that moving from 0 to 1 is interpreted as a 100\% increase. Results remain robust under this specification. 

Also connected to this are the event study results for the levels variables, which show a significant short-term increase followed by a sharp decline in the later periods. Two mechanisms could explain this trend. First, as illustrated in Figure \ref{fig:event_view_300km}, the group of municipalities treated at the earliest stage could have experienced a structural change in the final periods that drove levels downward. The point estimates are highly noisy and largely influenced by these municipalities. At the moment I have not found anything that can explain this, but additional evidence must be studied to discard this. Second, it is possible that municipalities with continuous case activity are driving the observed rise and fall in caseloads. However, this explanation is less convincing, since the extensive margin shows no effect and the same pattern does not appear across all outcomes, particularly those unrelated to prison outcomes. In general, further investigation need to be done, as the pattern is concerning and might be showing a bias on the overall results.

Finally, one additional point is worth noting. In the main results, I raised the possibility that estimates using rate-based dependent variables (e.g., the share of guilty individuals sentenced to prison or the ratio of pretrial detentions to individuals processed) could be biased. This arises because imputing zeros for municipalities with a denominator of zero requires adding a control: a dummy variable equal to one when the denominator variable is greater than zero and zero otherwise. The concern is that if treatment also affects the extensive margin, this dummy becomes endogenous to treatment and functions as a bad control. Figure \ref{fig:att_sent_robustness_chen}, however, shows that treatment is not correlated with this extensive-margin dummy for total processed and sentenced cases, alleviating this concern.

\subsubsection*{Robustness to different thresholds}

As specified in Section \ref{sec:emp}, treatment is defined using the minimum distance from each municipality to a federal prison. I adopt a data-driven thresholding approach: the median and mean of this distance are 236 km and 252 km, respectively, so I set the cutoff at 300 km—a round figure and a multiple of 100. This choice is further motivated by anecdotal evidence suggesting that federal prisons were intended to influence broader geographic zones rather than individual municipalities or entire states. Nevertheless, the choice of 300 km remains to some extent arbitrary. To assess the robustness of the results to alternative definitions of treatment, I re-estimate the same specification for the same outcomes using thresholds that vary by increments of +1 km within a +-50 km window around 300 km. Figure \ref{fig:att_sent_robust_thres} shows that both the point estimates and their associated standard errors remain nearly unchanged across these 100 alternative treatment definitions, providing strong evidence that the findings are robust to threshold selection. 

\subsubsection*{ITT and continuous treatment (discussion)}

One institutional feature that motivates the treatment definition is that, under Mexican law, individuals from municipality $x$ are typically assigned to serve their sentence in the nearest federal prison (with the exception of gang-related crimes, which are excluded from the analysis). However, because I generally do not observe the actual prison to which each individual was assigned, the estimated effect might not be the average treatment effect on the treated (ATT), which is what generally is estimated using DiD. Rather, it more closely resembles an intent-to-treat (ITT) effect: what I identify is the effect of being in a municipality located within the area covered by a federal prison, regardless of whether all individuals from that municipality were actually sent there. However, since the analysis is aggregated at the municipality level, an argument could be made to interpret the effects as increases in the probability of being sentenced to prison, given that your municipality is on the vicinity of a newly constructed prison.

A further assumption, related to this and to the previous section, is that treatment is binary—once a prison is constructed within a given threshold distance of a municipality I assign 1, 0 otherwise. This definition rests on two implicit assumptions. First, once a prison is built within the defined vicinity, municipalities just inside the threshold are assumed to be affected in the same way as those located directly adjacent to the prison. This assumption is plausible in light of the institutional context, and is reinforced by the argument that prison construction alleviates overcrowding at the state level; in practice, zones rather than individual municipalities are the affected units, this affected units are determined by distance to the center rather than borders (states) or demarcations, so the precise distance within the zone may matter less once the threshold is crossed.

The second assumption is that municipalities are not additionally affected if a second or third prison is constructed closer than the first. This assumption is less plausible, since the capacity channel would suggest that successive increases in capacity should yield additional effects. Nonetheless, for tractability, treatment is modeled as binary rather than continuous.

Ideally, treatment would instead be defined continuously, capturing the effect of reducing the distance to the nearest federal prison. This motivates the following alternative specification:

\begin{equation}
\label{eq:eq4}
    Y_{pt} = \phi + \beta \, DistanceInv_{pt} + \gamma_t + \tau_p + \epsilon_{pt},
\end{equation}

where $DistanceInv_{pt}$ is defined as
\begin{equation}
\label{eq:eq5}
DistanceInv_{pt} =
\begin{cases}
0, & \text{if } MinDistance_{pt} > Threshold, \\
\dfrac{MinDistance_{pt} - Threshold}{Threshold}, & \text{if } MinDistance_{pt} \leq Threshold,
\end{cases}
\end{equation}

with $MinDistance_{pt}$ denoting the minimum distance from municipality $p$ to the closest federal prison at time $t$. Recent literature on generalized difference-in-differences (DiD) designs has highlighted some challenges in correctly identifying unbiased effects when using a continuous treatment \citep{Cont-Chaise-AEA, CallawayGoodmanBaconSantAnna2024}. First and foremost, for correct interpretation of the effects, a stronger parallel trends assumption has to hold, which is that the trajectory of outcomes for lower-dose units must represent how the outcomes of higher-dose units would have evolved had they been exposed to the lower dose instead \citep{CallawayGoodmanBaconSantAnna2024}.

Unfortunately, none of these estimators can be applied to the context in hand. \cite{CallawayGoodmanBaconSantAnna2024} requires that the continuous treatment is absorbing, meaning that once it goes from 0 to a dose $d>0$, it has to stay in that level till the end of the study period. In my context, the continuous variable would increase once the distance to a federal prison decreases over time (when new prisons are built closer). \cite{Cont-Chaise-AEA} assumes that there are no stayers, meaning that all units change from period $t-1$ to period $t$ to a different dose. In my case, some units will remain at the same distance to a prison at different points in time. 

One important caveat is that there are concerns with using this specification, even if it were unbiased. First, constructing the treatment as a continuous variable would make the resulting coefficient difficult to interpret. At best, the interpretation would reduce to a relative measure of proximity on a 0–1 scale: municipalities with values closer to 1 are nearer to a federal prison. This interpretability issue does not arise with the current binary specification. Moreover, adopting a continuous treatment would make the parallel trends assumption considerably more difficult to satisfy, as it would need to hold across the group with low doses and the groups with high, rather than only between treated and untreated groups. Even with these concerns, adding this analysis is considered for future steps.

